\documentclass[../../main.tex]{subfiles}

\begin{document}

\section{外测度的构造}\label{section:ch17-outer-measure-construction}

我们通过首先定义赋予有界区间以它的长度值的原始集函数来构造实直线的子集上的 Lebesgue 外测度. 接着定义集合的外测度为覆盖这个集合的可数个有界区间的长度之和的下确界. 外测度的这种构造法在一般情形下也奏效.

\begin{theorem}\label{theorem:ch17-induced-outer-measure}
令 $\mathcal{S}$ 为集合 $X$ 的子集族而 $\mu:\mathcal{S}\to[0,\infty]$ 是一个集函数. 定义 $\mu^*(\varnothing)=0$, 且对 $E\subseteq X$, $E\neq\varnothing$ 定义
\begin{align}
\mu^*(E)=\inf\sum_{k=1}^{\infty}\mu(E_{k}),\label{align:ch17-outer-def}
\end{align}
其中下确界对所有覆盖 $E$ 的 $\mathcal{S}$ 中的集合的可数族 $\{E_{k}\}_{k=1}^{\infty}$ 取\footnote{我们允许 $E$ 的有限覆盖, 在这种情形下\eqref{align:ch17-outer-def}的和式是有限的. 若 $X$ 的子集 $E$ 不能被 $\mathcal{S}$ 的可数子族覆盖, 根据定义, 它有等于 $\infty$ 的外测度.}. 则集函数 $\mu^*:2^{X}\to[0,\infty]$ 是外测度, 称为由 $\mu$ 诱导的外测度.
\end{theorem}

\begin{proof}
为验证可数单调性, 令 $\{E_{k}\}_{k=1}^{\infty}$ 为覆盖集合 $E$ 的 $X$ 的子集族. 若对某个 $k$, $\mu^*(E_{k})=\infty$, 则 $\mu^*(E_{k})\leqslant\sum_{k=1}^{\infty}\mu^*(E_{k})=\infty$. 因此我们可以假设每个 $E_{k}$ 具有有限外测度. 令 $\varepsilon>0$. 对每个 $k$, 存在 $\mathcal{S}$ 中集合的可数族 $\{E_{ik}\}_{i=1}^{\infty}$, 它覆盖 $E_{k}$ 且
\begin{align*}
\sum_{i=1}^{\infty}\mu(E_{ik})<\mu^*(E_{k})+\frac{\varepsilon}{2^{k}}.
\end{align*}
则 $\{E_{ik}\}_{1\leqslant k,i<\infty}$ 是 $\mathcal{S}$ 中集合的可数族, 它覆盖 $\bigcup_{k=1}^{\infty}E_{k}$, 因此也覆盖 $E$. 由外测度的定义,
\begin{align*}
\mu^*(E)\leqslant\sum_{1\leqslant k,i<\infty}\mu(E_{ik})=\sum_{k=1}^{\infty}\left[\sum_{i=1}^{\infty}\mu(E_{ik})\right]\leqslant\sum_{k=1}^{\infty}\mu^*(E_{k})+\sum_{k=1}^{\infty}\varepsilon/2^{k}=\sum_{k=1}^{\infty}\mu^*(E_{k})+\varepsilon.
\end{align*}
由于该不等式对所有成立 $\varepsilon>0$, 它对 $\varepsilon=0$ 也成立.
\end{proof}

\begin{definition}[诱导的 Carath\'eodory 测度]\label{definition:ch17-induced-measure}
令 $\mathcal{S}$ 为 $X$ 的子集族, $\mu:\mathcal{S}\to[0,\infty]$ 是一个集函数, 而 $\mu^*$ 为由 $\mu$ 诱导的外测度. $\mu^*$ 在 $\mu^*$ 可测集的 $\sigma$ 代数 $\mathcal{M}$ 上的限制得到的测度 $\overline{\mu}$ 称为由 $\mu$ 诱导的 Carath\'eodory 测度.
\end{definition}

\begin{center}
\begin{tabular}{ccc}
$\mu^*:2^{X}\to[0,\infty]$ & & \\
(诱导的外测度) & & \\
$\uparrow$ & $\nwarrow$ & \\
& & \\
$\mu:\mathcal{S}\to[0,\infty]$ & & $\overline{\mu}:\mathcal{M}\to[0,\infty]$ \\
(一般的集函数) & & (诱导的 Carath\'eodory 测度) \\
\end{tabular}\\
{\it Carath\'eodory 构造法}
\end{center}

对于 $X$ 的子集族 $\mathcal{S}$, 我们用 $\mathcal{S}_{\sigma}$ 表示 $\mathcal{S}$ 中集合的可数并集, 而用 $\mathcal{S}_{\delta}$ 表示 $\mathcal{S}_{\sigma}$ 中集合的可数交集. 观察到若 $\mathcal{S}$ 是实数的开区间族, 则 $\mathcal{S}_{\sigma}$ 是 $\mathbf{R}$ 的开子集族而 $\mathcal{S}_{\delta}$ 是 $\mathbf{R}$ 的 $G_{\delta}$ 子集族.

我们证明了实数集 $E$ 是 Lebesgue 可测的当且仅当存在 $G_{\delta}$ 集 $G$ 使得 $G\sim E$ 的 Lebesgue 测度为零, 见 2.4 节定理 11. 以下命题告诉我们由一般的集函数诱导的 Carath\'eodory 测度的一个相关性质. 该性质是证明一些重要定理的关键要素, 其中有我们下一节要证明的 Carath\'eodory-Hahn 定理以及即将介绍的 Fubini 与 Tonelli 定理.

\begin{proposition}\label{proposition:ch17-gdelta-approximation}
令 $\mu:\mathcal{S}\to[0,\infty]$ 为定义在集合 $X$ 的子集族 $\mathcal{S}$ 上的集函数, 而 $\overline{\mu}:\mathcal{M}\to[0,\infty]$ 为由 $\mu$ 诱导的 Carath\'eodory 测度. 令 $E$ 为 $X$ 的子集, 满足 $\mu^*(E)<\infty$. 则存在 $X$ 的子集 $A$ 使得
\begin{align*}
A\in\mathcal{S}_{\sigma\delta},\quad E\subseteq A,\quad \mu^*(E)=\mu^*(A).
\end{align*}
此外, 若 $E$ 和 $\mathcal{S}$ 中的每个集关于 $\mu^*$ 可测, 则 $A$ 也如此, 且
\begin{align*}
\overline{\mu}(A\sim E)=0.
\end{align*}
\end{proposition}

\begin{proof}
令 $\varepsilon>0$. 我们宣称存在集 $A_{\varepsilon}$ 满足
\begin{align}
A_{\varepsilon}\in\mathcal{S}_{\sigma},\quad E\subseteq A_{\varepsilon}\text{ 且 }\mu^*(A_{\varepsilon})<\mu^*(E)+\varepsilon.\label{align:ch17-approx-eps}
\end{align}
事实上, 由于 $\mu^*(E)<\infty$, 存在覆盖 $E$ 的 $\mathcal{S}$ 中的集族 $\{E_{k}\}_{k=1}^{\infty}$ 使得
\begin{align*}
\sum_{k=1}^{\infty}\mu(E_{k})<\mu^*(E)+\varepsilon.
\end{align*}
定义 $A_{\varepsilon}=\bigcup_{k=1}^{\infty}E_{k}$. 则 $A_{\varepsilon}$ 属于 $\mathcal{S}_{\sigma}$ 且 $E\subseteq A_{\varepsilon}$. 此外, 由于 $\{E_{k}\}_{k=1}^{\infty}$ 是覆盖 $A_{\varepsilon}$ 的 $\mathcal{S}$ 中的集合的可数族, 根据外测度 $\mu^*$ 的定义,
\begin{align*}
\mu^*(A_{\varepsilon})\leqslant\sum_{k=1}^{\infty}\mu(E_{k})<\mu^*(E)+\varepsilon.
\end{align*}
因此\eqref{align:ch17-approx-eps}对 $A_{\varepsilon}$ 的这个选取成立.

定义 $A=\bigcap_{k=1}^{\infty}A_{1/k}$. 则 $A$ 属于 $\mathcal{S}_{\sigma\delta}$ 且 $E$ 是 $A$ 的子集, 因为 $E$ 是每个 $A_{1/k}$ 的子集. 此外, 根据 $\mu^*$ 的单调性和估计\eqref{align:ch17-approx-eps}, 对所有 $k$,
\begin{align*}
\mu^*(E)\leqslant\mu^*(A)\leqslant\mu^*(A_{1/k})\leqslant\mu^*(E)+\frac{1}{k}.
\end{align*}
因此 $\mu^*(E)=\mu^*(A)$.

现在假设 $E$ 是 $\mu^*$ 可测的且 $\mathcal{S}$ 中的每个集合是 $\mu^*$ 可测的. 由于可测集是 $\sigma$ 代数, 集合 $A$ 是可测的. 但 $\mu^*$ 是测度 $\overline{\mu}$ 的延拓. 因此, 根据测度的分割性质,
\begin{align*}
\overline{\mu}(A\sim E)=\overline{\mu}(A)-\overline{\mu}(E)=\mu^*(A)-\mu^*(E)=0.
\end{align*}
\end{proof}

\begin{exercise}\label{exercise:ch17-18}
令 $\mu^*:2^{X}\to[0,\infty]$ 为外测度. 令 $A\subseteq X$, $\{E_{k}\}_{k=1}^{\infty}$ 是可测集的不交可数族, 而 $E=\bigcup_{k=1}^{\infty}E_{k}$. 证明
\begin{align*}
\mu^*(A\cap E)=\sum_{k=1}^{\infty}\mu^*(A\cap E_{k}).
\end{align*}
\end{exercise}

\begin{exercise}\label{exercise:ch17-19}
证明任何由外测度诱导的测度是完备的.
\end{exercise}

\begin{exercise}\label{exercise:ch17-20}
令 $X$ 为任何集. 定义 $\eta:2^{X}\to[0,\infty]$, $\eta(\varnothing)=0$, 且对 $E\subseteq X$, $E\neq\varnothing$ 定义 $\eta(E)=\infty$. 证明 $\eta$ 是一个外测度. 证明赋予 $X$ 的每个子集 $0$ 的集函数是一个外测度.
\end{exercise}

\begin{exercise}\label{exercise:ch17-21}
令 $X$ 为一个集合, $\mathcal{S}=\{\varnothing,X\}$, 且定义 $\mu(\varnothing)=0$, $\mu(X)=1$. 确定由集函数 $\mu:\mathcal{S}\to[0,\infty]$ 诱导的外测度 $\mu^*$ 以及可测集的 $\sigma$ 代数.
\end{exercise}

\begin{exercise}\label{exercise:ch17-22}
在 $\mathbf{R}$ 的子集族 $\mathcal{S}=\{\varnothing,[1,2]\}$ 上定义集函数 $\mu:\mathcal{S}\to[0,\infty]$ 如下: $\mu(\varnothing)=0$, $\mu([1,2])=1$. 确定由 $\mu$ 诱导的外测度 $\mu^*$ 与可测集的 $\sigma$ 代数.
\end{exercise}

\begin{exercise}\label{exercise:ch17-23}
在 $\mathbf{R}$ 的所有子集构成的族 $\mathcal{S}$ 上, 通过设 $\mu(A)$ 为 $A$ 中的整数个数定义集函数 $\mu:\mathcal{S}\to\mathbf{R}$. 确定由 $\mu$ 诱导的外测度 $\mu^*$ 以及可测集的 $\sigma$ 代数.
\end{exercise}

\begin{exercise}\label{exercise:ch17-24}
令 $\mathcal{S}$ 为 $X$ 的子集族而 $\mu:\mathcal{S}\to[0,\infty]$ 是一个集函数. $\mathcal{S}$ 中的每个集合关于由 $\mu$ 诱导的外测度是可测的吗?
\end{exercise}

\end{document}